\chapter{VFS：super\_block、inode 与 dentry}
\label{chap:vfs}
miniOS 的 T113 VFS 使用 \code{super\_block}、\code{dentry}、\code{inode}、\code{file} 和\code{vfs\_mount} 组织挂载与打开文件。当前版本已经删除旧的\code{vfs\_backend\_ops} 和通用适配器 \code{backend\_fops}：每个文件系统直接提供自己的 \code{inode\_operations} 与 \code{file\_operations}，字符设备在打开时把\code{file->f\_op} 切换成驱动的操作表。

\section{核心对象}

\begin{ccode}
struct super_block {
    uint32_t s_magic;
    const struct super_operations *s_op;
    struct dentry *s_root;
    void *s_fs_info;
};

struct inode {
    uint32_t i_ino;
    uint32_t i_size;
    uint32_t i_type;
    uint32_t i_rdev;
    struct super_block *i_sb;
    const struct inode_operations *i_op;
    const struct file_operations *i_fop;
    void *i_private;
};

struct file {
    uint32_t f_flags;
    uint32_t f_pos;
    struct dentry *f_dentry;
    struct inode *f_inode;
    const struct file_operations *f_op;
    void *private_data;
    struct mutex f_pos_lock;
};
\end{ccode}

\begin{itemize}
  \item \code{super\_block} 表示一次文件系统挂载。
  \item \code{dentry} 表示路径中的名字，并指向 inode。
  \item \code{inode} 表示文件、目录或设备节点本身。
  \item \code{file} 表示一次 \code{open()}；\code{dup()} 和 \code{fork()} 可以共享它。
  \item \code{i\_rdev} 仅对字符设备有意义，保存 major/minor 合成的设备号。
  \item \code{private\_data} 给设备驱动保存每次打开的私有状态。
\end{itemize}

\section{三张操作表}

\begin{ccode}
struct super_operations {
    void (*put_super)(struct super_block *sb);
    void (*evict_inode)(struct inode *inode);
};

struct inode_operations {
    int (*lookup)(struct inode *dir, struct dentry *child);
    int (*create)(struct inode *dir, struct dentry *child);
    int (*truncate)(struct inode *inode);
};

struct file_operations {
    int (*open)(struct inode *inode, struct file *file);
    int (*read)(struct file *file, void *buf, uint32_t count);
    int (*write)(struct file *file, const void *buf, uint32_t count);
    int (*readdir)(struct inode *inode, const char *path,
                   uint32_t index, struct vfs_dirent *entry);
    void (*release)(struct inode *inode, struct file *file);
};
\end{ccode}

职责划分：

\begin{txtcode}
inode->i_op   路径和元数据：lookup/create/truncate
inode->i_fop  此 inode 默认的打开后操作
file->f_op    本次 open 实际使用的操作，可在 open 中替换
\end{txtcode}

\section{普通文件}

minifs、FAT32 和 bootfs 都直接提供自己的操作表：

\begin{txtcode}
minifs inode
+-- i_op  = minifs_iops
+-- i_fop = minifs_file_fops 或 minifs_dir_fops

FAT32 inode
+-- i_op  = fat32_iops
+-- i_fop = fat32_file_fops 或 fat32_dir_fops
\end{txtcode}

读取 \file{/bin/ls}：

\begin{txtcode}
sys_read(fd)
  → current->files[fd]
  → file->f_op->read(file, ...)
  → minifs_vfs_read()
  → 根据 file->f_inode->i_ino 读取 SD 卡 minifs 数据块
\end{txtcode}

读取 \file{/boot/kernel\_minios.img}：

\begin{txtcode}
sys_read(fd)
  → file->f_op->read
  → fat32_read()
  → file->f_inode->i_private（fat_inode）
  → FAT 簇链与 SD 块设备
\end{txtcode}

不再经过 \code{backend\_read()} 或 \code{super\_block->s\_backend}。

\section{字符设备}

devfs 只维护“名字到设备号”的关系：

\begin{txtcode}
"console" → inode(type=CHAR, i_rdev=1:0)
"ttyS0"   → inode(type=CHAR, i_rdev=3:0)
"key0"    → inode(type=CHAR, i_rdev=13:0)
\end{txtcode}

字符设备 inode 的默认操作是 \code{def\_chr\_fops}：

\begin{txtcode}
inode->i_fop = def_chr_fops
                   |
                   +-- open = chrdev_open
\end{txtcode}

\code{chrdev\_open()} 根据 \code{inode->i\_rdev} 查询全局字符设备表，然后替换本次打开文件的操作表：

\begin{ccode}
file->f_op = driver_fops;
\end{ccode}

因此 \file{/dev/console} 的完整路径是：

\begin{txtcode}
open("/dev/console")
  → devfs_iops.lookup
  → inode(i_rdev=1:0, i_fop=def_chr_fops)
  → file->f_op = inode->i_fop
  → def_chr_fops.open
  → chrdev_find(1:0)
  → file->f_op = console_fops

read(fd)
  → file->f_op->read
  → console_read
  → t113_uart_read
\end{txtcode}

GPIO 按键同样执行：

\begin{txtcode}
/dev/key0 → 13:0 → gpio_keys_fops → gpio_keys_read
\end{txtcode}

devfs 不再提供 \code{devfs\_read()} 或 \code{devfs\_write()}，也不接触 UART/GPIO寄存器。

\section{路径查找}

\code{lookup\_path()} 先以最长前缀选择挂载，再从该挂载的根 dentry 开始逐分量查找：

\begin{txtcode}
/bin/ls
  → minifs root inode
  → i_op->lookup(parent inode, child dentry "bin")
  → bin inode
  → i_op->lookup(bin inode, child dentry "ls")
  → ls inode
\end{txtcode}

文件系统的 \code{lookup()} 调用 \code{t113\_vfs\_bind\_inode()}，把查找结果绑定到dentry，并设置该 inode 的 \code{i\_op/i\_fop/i\_private}。inode cache 使用\code{(super\_block, i\_ino)} 作为键。

目前 devfs 直接使用 \code{child->d\_name}；minifs 和 FAT32 虽然已经原生实现\code{i\_op->lookup}，内部仍利用 \code{child->d\_path} 复用现有的完整路径解析器。后续可以把这两个解析器继续收窄为“父 inode + 单个分量”。

\section{挂载}

\begin{ccode}
int t113_vfs_mount_fs(const char *path,
                      uint32_t magic,
                      void *fs_info,
                      const struct super_operations *s_op,
                      const struct inode_operations *root_iop,
                      const struct file_operations *root_fop);
\end{ccode}

当前挂载：

\begin{txtcode}
/       minifs，SD 第二分区，只读
/boot   FAT32，SD 第一分区，读写；失败时可退回 bootfs
/dev    devfs，动态字符设备节点
/sys    sysfs，只读投影 driver core 的 kobject/device/bus/driver
\end{txtcode}

挂载根 inode 直接使用文件系统传入的 \code{root\_iop/root\_fop}。FAT32 通过\code{super\_operations.evict\_inode} 在最后一个 VFS inode 引用释放时写回目录项并释放 \code{fat\_inode} 引用。

\section{与 Linux 的关系}

当前主要派发方式已经与 Linux 一致：

\begin{txtcode}
路径操作       inode->i_op
普通文件 I/O   file->f_op → 文件系统 fops
字符设备 I/O   file->f_op → 驱动 fops
设备选择       inode->i_rdev → chrdev table
\end{txtcode}

miniOS 仍然比 Linux 简化：对象来自静态池，没有 RCU、页缓存、完整权限模型、负 dentry、umount 和并发 rename；挂载切换仍先用最长字符串前缀选挂载，而不是在逐级 path walk 中穿越 mountpoint。

\section{\code{vfs.h} 当前保留的重要结构体}

\code{vfs.h} 是文件系统、devfs 和设备驱动共同依赖的公开接口。当前重要结构体可以分成“操作表”“VFS 对象”和“目录项输出”三组。

\subsection{操作表}

\begin{txtcode}
super_operations
+-- put_super(sb)       卸载或释放文件系统实例（目前尚未实现 umount）
+-- evict_inode(inode)  VFS inode 最后一个引用消失时释放文件系统私有对象

inode_operations
+-- lookup(dir, child)  在目录中查找一个名字，并给 child 绑定 inode
+-- create(dir, child)  创建普通文件并绑定 inode
+-- truncate(inode)     截断文件

file_operations
+-- open(inode, file)   初始化本次打开；也可以替换 file->f_op
+-- read(file, ...)
+-- write(file, ...)
+-- readdir(inode, path, index, entry)
+-- release(inode, file)
\end{txtcode}

\code{inode\_operations} 操作的是文件系统名字空间和文件元数据；\code{file\_operations} 操作的是一次已经打开的文件。当前 \code{readdir()} 仍以 inode而不是打开的目录 file 为参数，这是 miniOS 的简化，并非完整 Linux 接口。

\subsection{VFS 对象}

\begin{txtcode}
vfs_mount
+-- path                 挂载路径，例如 /、/boot、/dev
+-- sb -------------------------------+
+-- mnt_root ----------------+        |
+-- parent                   |        |
+-- mountpoint               |        |
                             v        v
                         dentry    super_block
                         +-- d_name   +-- s_magic
                         +-- d_path   +-- s_op
                         +-- d_parent +-- s_root --> 根 dentry
                         +-- d_sb ----+
                         +-- d_inode  +-- s_fs_info
                                |
                                v
                              inode
                              +-- i_ino / i_size / i_type
                              +-- i_rdev
                              +-- i_sb
                              +-- i_op
                              +-- i_fop
                              +-- i_private

进程 fd 表 --> file
               +-- f_dentry --> dentry
               +-- f_inode  --> inode
               +-- f_op
               +-- f_pos / f_pos_lock
               +-- private_data
\end{txtcode}

各对象分别回答不同问题：

\begin{table}[H]
\centering
\small
\begin{tabularx}{\linewidth}{l>{\raggedright\arraybackslash}X>{\raggedright\arraybackslash}X}
\toprule
对象 & 表示什么 & 关键关系 \\
\midrule
\code{super\_block} & 一次文件系统实例 & \code{s\_root} 指向根 dentry；\code{s\_fs\_info} 保存 minifs/FAT32/devfs 的实例数据 \\
\code{dentry} & 路径中的一个名字 & \code{d\_parent} 构成目录层次，\code{d\_inode} 指向名字对应的 inode \\
\code{inode} & 文件、目录或设备节点本身 & \code{i\_op} 管路径/元数据，\code{i\_fop} 是默认打开操作 \\
\code{file} & 一次 \code{open()} 形成的打开文件描述 & \code{f\_op} 是本次打开最终使用的操作表；\code{f\_pos} 是共享偏移 \\
\code{vfs\_mount} & 把一个 super block 接到路径树上 & 保存挂载路径、父挂载和挂载点 \\
\code{vfs\_dirent} & \code{readdir()} 返回给调用者的一项 & 只是结果结构，不参与缓存和引用关系 \\
\bottomrule
\end{tabularx}
\end{table}

三个私有指针不要混淆：

\begin{txtcode}
super_block.s_fs_info  每次挂载/文件系统实例的私有数据
inode.i_private       某个 inode 的文件系统私有数据，例如 FAT inode
file.private_data     某次 open 的私有数据，例如设备会话状态
\end{txtcode}

\subsection{引用和共享关系}

当前对象来自 \code{vfs.c} 的固定大小静态池：8 个 mount、8 个 super block、128 个inode、128 个 dentry 和 32 个 file。它们不是动态分配的。

\begin{txtcode}
mount 持有 super_block 和根 dentry
dentry 持有 inode，并持有父 dentry
file 持有 dentry；通过 dentry 间接保持 inode 有效
进程 files[fd] 持有 file
dup()/fork() 增加同一个 file->references
\end{txtcode}

因此 \code{dup()} 或 \code{fork()} 后多个 fd/进程共享同一个 \code{struct file}，也共享\code{f\_pos}；普通文件用 \code{f\_pos\_lock} 串行化 I/O 和偏移更新。字符设备不使用文件位置，所以 \code{t113\_vfs\_read/write()} 不为字符设备获取 \code{f\_pos\_lock}。

\code{vfs\_lock} 只保护对象池、引用计数和 file 有效性检查。读写时 VFS 先在锁内临时增加 \code{file->references}，随后释放全局锁，再调用 \code{file->f\_op->read/write()}；这样 UART 阻塞读取时不会长期占用全局 VFS 锁。

\section{\code{chrdev.h} 当前保留的内容}

\code{chrdev.h} 刻意保持很薄，目前没有公开 \code{struct cdev} 或\code{struct chrdev\_entry}。它只包含：

\begin{ccode}
MINIOS_MKDEV(major, minor)  /* 组合 16 位 major 和 16 位 minor */
MINIOS_MAJOR(dev)
MINIOS_MINOR(dev)

void t113_chrdev_init(void);
int t113_chrdev_register(uint16_t major, uint16_t minor,
                         const struct file_operations *fops);
const struct file_operations *t113_chrdev_inode_fops(void);
\end{ccode}

真正的注册表结构是 \code{chrdev.c} 的私有实现：

\begin{ccode}
struct chrdev_entry {
    uint32_t dev;                         /* 完整 major:minor */
    const struct file_operations *fops;  /* 具体设备驱动操作 */
    uint32_t used;
};

static struct chrdev_entry chrdev_table[32];
\end{ccode}

这种封装意味着 VFS、devfs 和驱动都不能直接遍历或修改全局字符设备表，只能调用注册接口。\code{chrdev\_lock} 保护注册和查找。当前没有注销接口、设备号范围、动态major、\code{struct cdev}、引用计数或驱动卸载，因此它更接近 Linux 2.4 的简单注册表，而不是 Linux 2.6 以后的 \code{cdev\_map}。

\section{\code{vfs.h} 与 \code{chrdev.h} 的关系}

\code{chrdev.h} 依赖 \code{vfs.h}，因为字符设备最终注册的是 VFS 定义的\code{struct file\_operations}：

\begin{txtcode}
具体驱动
  +-- console_fops / gpio_keys_fops
          |
          | t113_chrdev_register(major, minor, fops)
          v
chrdev_table：dev_t -> driver fops

devfs
  +-- 名字 -> dev_t
      例如 "console" -> 1:0

VFS inode
  +-- i_type = VFS_NODE_CHAR
  +-- i_rdev = 1:0
  +-- i_fop = def_chr_fops
                    | open
                    v
              chrdev_open()
                    | 根据 i_rdev 查 chrdev_table
                    v
              file->f_op = console_fops
\end{txtcode}

这里实际存在两张彼此独立的表：

\begin{table}[H]
\centering
\small
\begin{tabularx}{\linewidth}{>{\raggedright\arraybackslash}X>{\raggedright\arraybackslash}X>{\raggedright\arraybackslash}X}
\toprule
表 & 保存的映射 & 负责什么 \\
\midrule
devfs \code{nodes[]} & \code{"console" -> (1,0)} & 路径名字和设备号的关系 \\
chrdev \code{chrdev\_table[]} & \code{(1,0) -> console\_fops} & 设备号和驱动操作的关系 \\
\bottomrule
\end{tabularx}
\end{table}

\code{t113\_devfs\_create\_chrdev("console", 1, 0)} 只创建第一种映射；\code{t113\_chrdev\_register(1, 0, \&console\_fops)} 只创建第二种映射。两步都成功后，\file{/dev/console} 才能正常打开并到达 UART 驱动。

\section{从系统调用到对象的完整关系}

以 \code{read(fd, buffer, count)} 为例：

\begin{txtcode}
用户 fd（整数）
  -> current->files[fd]             struct file *
  -> t113_vfs_read(file, ...)
       +-- 临时增加 file 引用
       +-- 普通文件：锁住 file->f_pos_lock
       +-- file->f_op->read(file, ...)
              +-- 普通文件：minifs/FAT32 fops
              |      -> file->f_inode
              |      -> inode->i_private / inode->i_sb->s_fs_info
              |      -> SD 块设备
              +-- 字符设备：console_fops
                     -> UART RX 缓冲区/寄存器
\end{txtcode}

用户态的 fd 只是当前进程 fd 表的数组下标；内核 \code{struct file} 才是打开文件对象。它既不是磁盘上的文件，也不是 inode。多个 fd 可以指向同一个 file，一个inode 也可以同时对应多个相互独立、偏移不同的 file。

\section{sysfs：把 driver core 投影到 \file{/sys}}

目录如何生成、属性怎样由驱动提供、与 Linux sysfs 的差别，详见第~\ref{chap:sysfs} 章。

\code{CONFIG\_SYSFS=y} 时，\code{kernel\_boot()} 在 driver core 完成设备枚举与驱动绑定后调用 \code{t113\_sysfs\_init()}。根 minifs 只提供空的 \file{/sys} 挂载点，实际目录和属性由 \code{sysfs.c} 根据当前内存对象动态生成，不写入 SD 卡。

\begin{txtcode}
/sys
+-- devices                         按 device.parent 展开的设备拓扑
|   +-- platform
|       +-- gic
|       +-- uart3
|       +-- timer0
|       +-- user-key
|       +-- twi2
|       +-- sdc1
|           +-- mmc1:0001:1
+-- bus                             按 bus_type 分类的视图
    +-- platform
    |   +-- devices
    |   +-- drivers
    +-- i2c
    |   +-- devices
    |   +-- drivers
    +-- sdio
        +-- devices
        +-- drivers
\end{txtcode}

设备目录公开 \code{name}、\code{bus}、\code{driver}、\code{path} 四个只读属性；驱动目录公开\code{name} 与 \code{bound}。例如：

\begin{shcode}
ls /sys/devices/platform
cat /sys/devices/platform/user-key/driver
cat /sys/devices/platform/user-key/value
cat /sys/devices/platform/user-key/keycode
cat /sys/bus/platform/drivers/gpio-keys/bound
cat /sys/bus/sdio/devices/mmc1:0001:1/path
\end{shcode}

当前 VFS 没有符号链接 inode，因此 \file{/sys/bus/*/devices/*} 使用可浏览的镜像目录，其 \code{path} 属性指回 \file{/devices/...}；这与 Linux sysfs 使用符号链接表达同一关系在语义上相同，但表示形式更简单。属性在每次 \code{read()} 时实时读取 driver core，后续注册的设备也会出现在目录中。

\section{GPIO 按键}

probe、中断、事件队列、寄存器和已知问题详见 第~\ref{chap:sysfs} 章。

100ASK T113-S3 RESET 旁的 USER KEY 已作为 platform device \code{user-key} 注册：

\begin{txtcode}
compatible  gpio-keys
GPIO        PB4
有效电平    低
上下拉      内部上拉
中断        GPIOB，双边沿
消抖        20 ms
键值        KEY_ENTER (28)
\end{txtcode}

\code{gpio\_keys} 驱动在 probe 中申请 GPIOB IRQ，建立 16 项事件环形队列和等待队列，并注册字符设备 \code{(13,0)} 与 devfs 节点 \file{/dev/key0}。按下和释放分别产生\code{value=1}、\code{value=0} 的 \code{input\_event}；队列为空时读取进程睡眠，由 IRQ 入队后唤醒。用户态可运行：

\begin{shcode}
keytest
\end{shcode}

驱动模型与设备 I/O 是两条互补路径：

\begin{txtcode}
/sys/devices/platform/user-key  查看设备、总线和绑定驱动
/dev/key0                       阻塞读取实际按键事件
\end{txtcode}

\section{已知限制（主机侧测试结果）}

把 \code{vfs.c}、minifs、FAT32、bio、devfs、chrdev 用主机 gcc（ASan/UBSan）编进一个测试程序，用 \code{make t113} 生成的 \code{minifs.img} 和一个 FAT32 分区拼成 SD 卡镜像，按 \code{kernel\_boot()} 的顺序挂载后测试。下面这些都通过：

\begin{itemize}
  \item 挂载 \file{/}、\file{/boot}、\file{/dev}；读 \file{/init}、\file{/bin/sh}、\file{/bin/init}（硬链接）、\file{/boot/sub/deep.txt}，内容与主机文件一致；
  \item \code{readdir} 列出 \file{/}、\file{/bin}、\file{/boot}、\file{/dev}；
  \item \file{/dev/console} 经 \code{chrdev\_open()} 切到驱动 fops，读写成功，\code{f\_pos} 保持 0；
  \item FAT32 的 \code{O\_CREATE}、写入、重新打开、\code{O\_TRUNC}；
  \item 同一文件打开两次共享一个 dentry 和 inode；
  \item 反复打开/关闭 500 次后 dentry、inode、file 槽位数回到开机时的值，FAT32 的16 个 \code{fat\_inode} 槽位也没有耗尽；没有锁重入。
\end{itemize}

发现的问题：

\begin{enumerate}
  \item \textbf{路径中间是非目录时没有拒绝。} \code{lookup\_path()} 不检查 \code{current} 是否为目录，直接调用它的 \code{i\_op->lookup}。minifs 和 FAT32 用完整路径解析，会自己失败；devfs 只比较 \code{child->d\_name}，所以 \file{/dev/console/console} 能打开控制台。在 \code{lookup\_path()} 调 \code{lookup} 之前加一次 \code{i\_type == VFS\_NODE\_DIR} 检查即可，这对应 Linux 的 \code{-ENOTDIR}。
  \item \textbf{只写的字符设备打不开。} \code{t113\_chrdev\_register()} 允许只有 \code{write} 的fops，但 \code{t113\_vfs\_open()} 总是要求 \code{f\_op->read}。现在只有 \code{RDONLY/RDWR}两种模式，没有 \code{O\_WRONLY}，所以这种设备任何模式都打不开。
  \item \code{alloc\_inode()} 命中缓存时传给 \code{evict\_inode()} 的临时 \code{struct inode} 只填了部分字段（\code{i\_op}、\code{i\_fop} 未初始化）。FAT32 只读 \code{i\_private}，现在没问题，但新文件系统的 \code{evict\_inode} 不能依赖其他字段。
  \item bootfs 的目录 fops 没有 \code{readdir}，退回 bootfs 时 \code{ls /boot} 会失败。
  \item 挂载根 inode 的 \code{i\_ino} 固定为 0，minifs 真正的根是 1；\file{/bin/..} 会得到另一个\code{i\_ino = 1} 的 inode。只读情况下没有影响。
\end{enumerate}
